\chapter{Quando la Giustizia Sbaglia: I Bias Cognitivi in Tribunale}

Nei capitoli precedenti abbiamo visto come il nostro cervello del Pleistocene si trovi spesso in difficoltà di fronte alle sfide del mondo moderno. Abbiamo esplorato come il Sistema 1 di Kahneman ,veloce, automatico e basato su euristiche ,possa portarci fuori strada quando si trova ad affrontare problemi complessi che richiederebbero l'intervento del più lento e riflessivo Sistema 2. Abbiamo visto come il bias di conferma ci faccia cercare informazioni che confermano le nostre convinzioni preesistenti, come l'effetto ancoraggio influenzi le nostre valutazioni economiche, e come l'euristica della disponibilità distorca la nostra percezione dei rischi.

Ma cosa succede quando questi stessi meccanismi cognitivi ,che erano perfettamente adatti per prendere decisioni rapide in piccoli gruppi tribali ,si manifestano nell'ambiente che dovrebbe essere il tempio della razionalità e dell'imparzialità: l'aula di tribunale?

La risposta è tanto semplice quanto inquietante: anche lì, sotto le toghe e i codici, battono cervelli umani con tutte le loro imperfezioni evolutive. E quando queste imperfezioni si scatenano in un contesto dove si decidono la libertà, la dignità e talvolta la vita delle persone, le conseguenze possono essere devastanti.

\section{L'illusione della razionalità giudiziaria}

Quando pensiamo a un giudice, l'immagine che ci viene in mente è quella di una figura saggia e imparziale, capace di separare i fatti dalle emozioni e di applicare la legge con rigore scientifico. È un'immagine rassicurante, costruita su secoli di tradizione giuridica che ha fatto della razionalità e dell'obiettività i suoi pilastri fondamentali.

Ma negli ultimi decenni, la ricerca scientifica ha iniziato a sollevare il velo su questa illusione. Nel 2001, un gruppo di ricercatori ,Chris Guthrie, Jeffrey Rachlinski e Andrew Wistrich ,condusse uno studio rivoluzionario coinvolgendo 167 giudici federali americani. L'obiettivo era semplice quanto ambizioso: testare se i magistrati fossero realmente immuni dai bias cognitivi che affliggono il resto dell'umanità\footnote{Chris Guthrie, Jeffrey J. Rachlinski e Andrew J. Wistrich, "Inside the Judicial Mind", Cornell Law Review, Vol. 86, 2001.}.

I risultati furono sconcertanti. I giudici si dimostrarono vulnerabili agli stessi errori sistematici di ragionamento che caratterizzano qualsiasi altro essere umano. Non perché fossero incompetenti o disonesti, ma semplicemente perché erano umani.

In Italia, studi condotti dall'Università di Milano-Bicocca in collaborazione con il Consiglio Superiore della Magistratura hanno rivelato pattern simili. Un'analisi di oltre 2.000 sentenze penali ha mostrato che anche i magistrati italiani sono soggetti agli stessi pregiudizi cognitivi documentati nella ricerca internazionale, con particolare evidenza per l'effetto àncora nelle determinazioni delle pene pecuniarie e l'euristica della disponibilità nei reati che ricevono maggiore copertura mediatica\footnote{Garofoli, R., \& Marchetti, D. (2019). "Bias cognitivi nelle decisioni giudiziarie: uno studio empirico sulla magistratura italiana". \textit{Cassazione Penale}, 59(4), 1523-1547.}.

Prendiamo l'effetto ancoraggio, che abbiamo già incontrato nel capitolo dedicato agli investimenti finanziari. In un esperimento, alcuni giudici dovevano valutare il risarcimento per un grave incidente automobilistico dove un pedone era rimasto paralizzato a vita. A un gruppo non venne fornita alcuna cifra di riferimento; a un altro venne detto che il pubblico ministero aveva chiesto l'archiviazione perché il danno non raggiungeva i 75.000 dollari.

Il risultato fu impressionante: i giudici senza "àncora" assegnarono un risarcimento medio di 1.249.000 dollari, mentre quelli esposti alla cifra di riferimento si fermarono a 882.000 dollari\footnote{Giuseppe Motta, "La rilevanza dei Bias cognitivi nell'attività giudiziaria", Laboratorio di sociologia del diritto, disponibile online.}. Una cifra completamente arbitraria, menzionata di sfuggita, aveva influenzato decisioni che avrebbero cambiato per sempre la vita delle persone coinvolte.

Ma c'è di più. Ricordate il bias del punto cieco di cui abbiamo parlato? Quel meccanismo per cui tutti crediamo di essere più razionali e meno influenzabili della media? Ebbene, uno studio ha rivelato che il 97\% dei giudici si considera più bravo della media nell'evitare i bias razziali ,una statistica matematicamente impossibile che dimostra come anche i magistrati cadano nella trappola dell'overconfidence\footnote{Laboratorio di sociologia del diritto, "Bias cognitivi: ovvero come i pregiudizi influiscono sul ragionamento", disponibile online.}.

È come se ogni giudice fosse convinto di essere l'unico lucido in un mondo di pazzi, senza rendersi conto di far parte dello stesso sistema che critica.

\section{Quando il Sistema 1 prende il controllo dell'aula}

Ricordate la distinzione tra Sistema 1 e Sistema 2 che abbiamo esplorato nei primi capitoli? Il Sistema 1 ,veloce, automatico, basato su pattern recognition ,e il Sistema 2 ,lento, deliberativo, che richiede sforzo cognitivo? Abbiamo visto come il Sistema 2 sia "pigro" e tenda a delegare al Sistema 1 ogni volta che può, specialmente quando è stanco o sotto pressione.

Questa dinamica si manifesta in modo drammatico nei tribunali. Studi hanno dimostrato che i giudici che hanno sentito molti casi in una giornata tendono a dare sentenze più severe verso sera. I magistrati mentalmente affaticati si affidano maggiormente a scorciatoie cognitive, stereotipi e prime impressioni\footnote{State of Mind, "Processo decisionale del giudice: aspetti cognitivi delle decisioni che fanno sentenza", 2017.}.

È come se, quando il "supervisore razionale" che abbiamo descritto nei capitoli precedenti è troppo stanco per fare il suo lavoro, l'aula di tribunale venisse lasciata in balia del pilota automatico cognitivo. Il giudice continua a credere di star prendendo decisioni ponderate e razionali, ma in realtà sta operando con gli stessi meccanismi mentali che usereste per scegliere cosa mangiare a colazione.

Il problema è che una colazione sbagliata vi rovina al massimo la giornata. Una sentenza sbagliata può rovinare una vita.

Ma qui entriamo in territorio ancora più insidioso. Come abbiamo visto parlando del "ragionamento motivato", il Sistema 2 spesso non corregge gli errori del Sistema 1, ma si limita a trovare giustificazioni razionali per decisioni già prese istintivamente. Il giudice stanco che ha una prima impressione negativa di un imputato non si ferma a riconsiderare quella impressione. Invece, il suo Sistema 2 inizia a cercare elementi nel fascicolo che giustifichino quella sensazione iniziale.

È la versione giudiziaria di quello che succede quando decidiamo istintivamente di comprare qualcosa e poi ci convinciamo che è un acquisto razionale elencando tutte le ragioni per cui ne avevamo "bisogno".

\section{L'arsenale completo dei bias giudiziari}

\subsection{Il bias di conferma: quando la verità diventa nemica}

Il bias di conferma, che abbiamo visto all'opera nei social media e nella polarizzazione politica, trova nel sistema giudiziario il suo terreno più letale. Una volta che investigatori, pubblici ministeri o giudici si sono fatti un'idea su un caso, tendono inconsciamente a cercare prove che confermino quella teoria iniziale, ignorando o minimizzando tutto ciò che potrebbe contraddirla.

Questo meccanismo, che nei nostri antenati garantiva coerenza nelle decisioni tribali, oggi può trasformarsi in una trappola mortale per l'amministrazione della giustizia. Il National Institute of Justice americano ha condotto un'analisi sistematica di 50 casi di condanne errate, identificando il bias di conferma come "il fattore causale più connesso con le condanne errate, per un margine significativo"\footnote{National Institute of Justice, "Confirmation Bias and Other Systemic Causes of Wrongful Convictions: A Sentinel Events Perspective", 2020.}.

Gli esperti lo chiamano "visione a tunnel": una volta che gli investigatori si convincono di aver identificato il colpevole, iniziano a interpretare ogni elemento come prova di quella colpevolezza. Le testimonianze ambigue diventano conferme, i comportamenti innocenti diventano sospetti, le coincidenze diventano prove schiaccianti. È come se indossassero lenti colorate che fanno apparire tutto dello stesso colore.

In Italia, oltre al caso Tortora che abbiamo già menzionato, il caso di Stefano Cucchi rappresenta un altro esempio tragico di visione a tunnel investigativa. Le indagini iniziali si concentrarono ossessivamente sull'ipotesi di morte per cause naturali o per uso di stupefacenti, ignorando sistematicamente le evidenze fisiche di percosse. Solo anni dopo, grazie all'ostinazione della famiglia e a nuove perizie mediche, emerse la verità sui pestaggi subiti in caserma\footnote{Anselmi, T., \& Gonnella, P. (2020). "Il caso Cucchi: anatomia di un depistaggio". \textit{Antigone - Quadrimestrale di critica del sistema penale}, 15(1), 78-95.}.

Il caso di Olindo Romano e Rosa Bazzi — condannati per la strage di Erba — continua a sollevare interrogativi sulla possibile influenza del pregiudizio di conferma, con alcuni elementi investigativi che sembrano essere stati interpretati forzatamente per adattarsi alla teoria accusatoria iniziale\footnote{Monza, G. (2021). "Dubbi sulla strage di Erba: tra certezze processuali e questioni irrisolte". \textit{Diritto Penale Contemporaneo}, 3, 112-134.}.

\subsection{L'euristica della disponibilità in tribunale}

L'euristica della disponibilità, che abbiamo visto influenzare le nostre percezioni dei rischi quotidiani, ha effetti particolarmente insidiosi nel sistema giudiziario. I giudici tendono a sovrastimare la probabilità di eventi che sono facilmente richiamabili alla memoria, spesso a causa della copertura mediatica.

Uno studio illuminante ha dimostrato che i giurati esposti a notizie su ingenti risarcimenti civili nei giorni precedenti al processo assegnavano danni mediamente sei volte superiori rispetto ai gruppi di controllo: 1.286.000 dollari contro una media di 96.000-226.000 dollari\footnote{The Jury Expert, "Media Exposure, Juror Decision-Making, and the Availability Heuristic", 2012.}. La copertura mediatica aveva letteralmente riprogrammato le loro aspettative su cosa fosse un risarcimento "appropriato".

È lo stesso meccanismo che ci fa sovrastimare il rischio di attentati terroristici dopo averli visti in televisione, o che ci convince che certi crimini siano più frequenti di quanto non siano realmente solo perché ne sentiamo parlare di più.

Pensateci: se negli ultimi mesi avete sentito molto parlare di un particolare tipo di reato sui media, questo influenzerà inconsciamente la vostra percezione di quanto sia "normale" che accada. E se fate parte di una giuria, questa percezione distorta potrebbe influenzare il vostro giudizio sulla colpevolezza o sull'entità della pena.

\subsection{Il tribalismo cognitivo nell'era della giustizia}

Ricordate il bias dell'appartenenza al gruppo che abbiamo analizzato nel capitolo sui social media? Quella tendenza automatica a favorire chi percepiamo come appartenente al nostro "gruppo" e a svalutare chi consideriamo "diverso"? Ebbene, questo meccanismo opera anche nelle aule di tribunale, spesso in modi sottili ma devastanti.

I dati dell'Innocence Project americano sono agghiaccianti: il 58\% delle persone scagionate da errori giudiziari è afroamericano, nonostante questa popolazione rappresenti solo il 13,6\% della popolazione totale. Gli afroamericani hanno una probabilità sette volte maggiore di essere erroneamente condannati per omicidio rispetto ai bianchi\footnote{Montana Innocence Project, "More than half of all exonerees are Black. Here's why", disponibile online.}.

Questo non accade perché giudici e giurati siano consciamente razzisti. Il problema è più sottile e insidioso: associazioni inconsce tra razza e criminalità, ereditate da secoli di storia e rinforzate da rappresentazioni mediatiche distorte, influenzano migliaia di micro-decisioni durante tutto il processo, dalla credibilità attribuita ai testimoni alla severità delle pene.

Una ricerca della Cornell University che ha analizzato migliaia di trascrizioni di selezioni di giurie ha scoperto che i pubblici ministeri escludevano i potenziali giurati afroamericani a un tasso oltre tre volte superiore rispetto ai bianchi nei casi di pena capitale\footnote{Cornell University Department of Statistics and Data Science, "Data Science analysis of court transcripts reveals biased jury selection", 2023.}. Il sistema che dovrebbe garantire un processo equo stava sistematicamente eliminando la diversità dalle giurie.

È il tribalismo cognitivo che abbiamo visto operare sui social media, trasportato nelle aule di tribunale dove può decidere della libertà delle persone.

\subsection{Il bias del senno di poi: quando il futuro colora il passato}

C'è un altro bias particolarmente insidioso nel contesto giudiziario: il bias del senno di poi, o hindsight bias. Una volta che sappiamo come è andata a finire una storia, tendiamo a vedere gli eventi che l'hanno preceduta come più prevedibili di quanto non fossero realmente.

Immaginate di dover giudicare se un medico ha commesso negligenza. Se il paziente è morto, tenderete a vedere ogni decisione del medico come "ovviamente" sbagliata, anche se in base alle informazioni disponibili al momento quella decisione era ragionevole. È difficile mettersi nei panni di chi doveva decidere senza sapere come sarebbe andata a finire.

Questo bias è particolarmente problematico nei casi dove si deve valutare la responsabilità per eventi drammatici. I giurati, sapendo già che è successo qualcosa di terribile, tendono a vedere i segnali premonitori ovunque e a giudicare più severamente chi "avrebbe dovuto prevedere" l'esito\footnote{Journal of Empirical Legal Studies, "Biases in legal decision-making: Comparing prosecutors, defense attorneys, law students, and laypersons", 2023.}.

\section{L'algoritmo che giudica: il caso COMPAS}

Se i bias umani erano già abbastanza preoccupanti, l'introduzione dell'intelligenza artificiale nel sistema giudiziario ha aperto scenari ancora più complessi. Negli Stati Uniti, migliaia di decisioni giudiziarie sono oggi influenzate da algoritmi di "valutazione del rischio" che promettono di rendere la giustizia più oggettiva e scientifica.

Il caso più famoso è quello di Eric Loomis, un cittadino afroamericano del Wisconsin che nel 2013 venne arrestato per aver guidato un'auto precedentemente utilizzata in una sparatoria. Pur non essendo coinvolto nella sparatoria stessa, patteggiò per due capi d'accusa minori: tentativo di fuga da un ufficiale e utilizzo di un veicolo senza consenso del proprietario.

Quando arrivò il momento di determinare la pena, il giudice utilizzò un sistema chiamato COMPAS (Correctional Offender Management Profiling for Alternative Sanctions), un algoritmo sviluppato dalla società privata Northpointe per valutare il rischio di recidiva. COMPAS analizza 137 fattori diversi ,dalla storia criminale alle condizioni socioeconomiche, dalle reti sociali ai pattern comportamentali ,per produrre un punteggio da 1 a 10 che indica la probabilità che una persona commetta nuovamente un reato\footnote{Stefano Carrer, "Se l'amicus curiae è un algoritmo: il chiacchierato caso Loomis alla Corte Suprema del Wisconsin", Giurisprudenza Penale Web, 2019.}.

COMPAS classificò Loomis come "ad alto rischio", e il giudice lo condannò a sei anni di carcere. Il problema era che COMPAS era una "scatola nera": l'algoritmo era coperto da segreto industriale, e nessuno ,nemmeno il giudice ,poteva sapere esattamente come arrivasse alle sue conclusioni.

Loomis fece appello, sostenendo che essere giudicato da un algoritmo segreto violava il suo diritto a un processo equo. La Corte Suprema del Wisconsin respinse l'appello nel 2016, stabilendo che l'uso di COMPAS era legittimo purché non fosse l'unico fattore determinante nella decisione\footnote{Supreme Court of Wisconsin, State of Wisconsin v. Eric L. Loomis, 13 luglio 2016.}.

Ma le investigazioni successive rivelarono un problema ancora più profondo. Nel 2016, un'inchiesta di ProPublica analizzò oltre 7.000 casi in cui era stato utilizzato COMPAS e scoprì che l'algoritmo aveva significativi bias razziali: aveva il doppio delle probabilità di classificare erroneamente gli imputati neri come "alto rischio" rispetto ai bianchi, mentre i bianchi avevano maggiori probabilità di essere erroneamente classificati come "basso rischio"\footnote{ProPublica, "Machine Bias: Risk Assessments in Criminal Sentencing", 2016.}.

L'ironia era perfetta: nel tentativo di eliminare i bias umani dalla giustizia, avevamo creato un sistema che codificava e amplificava quei bias in codice informatico. Come abbiamo visto nel capitolo sull'intelligenza artificiale, gli algoritmi non sono neutri: sono il riflesso digitale dei dati su cui vengono addestrati, e se quei dati riflettono secoli di discriminazioni storiche, l'algoritmo le perpetuerà e le renderà più efficienti.

È come avere un specchio magico che, invece di mostrarci la realtà, riflette e amplifica tutti i nostri pregiudizi, rendendoli però "oggettivi" perché vengono da un computer.

\subsection{Altri sistemi di valutazione: tra speranze e delusioni}

COMPAS non è l'unico sistema di questo tipo. Il Public Safety Assessment (PSA), sviluppato utilizzando dati di oltre 750.000 casi provenienti da circa 300 giurisdizioni, utilizza solo nove fattori tratti da registrazioni amministrative per valutare il rischio di fuga e di nuovi reati prima del processo. A differenza di COMPAS, il PSA è gratuito e i suoi metodi sono pubblicamente accessibili\footnote{Brookings Institution, "Understanding risk assessment instruments in criminal justice", 2019.}.

In teoria, il PSA dovrebbe essere più equo perché evita domande soggettive e si basa solo su dati "oggettivi" come l'età al momento dell'arresto e il numero di precedenti penali. Ma anche qui si nasconde un problema: come abbiamo visto, il numero di precedenti penali è fortemente correlato alla razza, non necessariamente perché alcune comunità commettano più reati, ma perché sono sottoposte a controlli più intensi da parte delle forze dell'ordine.

È come misurare l'altezza con un metro che ha segni diversi per persone diverse: il risultato sembrerà "oggettivo", ma sarà sistematicamente distorto.

Due ricercatori del Dartmouth College, Julia Dressel e Hany Farid, hanno condotto un esperimento illuminante: hanno confrontato l'accuratezza di COMPAS con quella di persone comuni a cui venivano forniti solo due dati sull'imputato: età e numero di precedenti penali. Il risultato? L'algoritmo sofisticato non era più accurato delle valutazioni umane basate su questi due semplici fattori\footnote{Aequitas Magazine, "Il caso Compas e il tema della giustizia predittiva, tra opacità e distorsioni", 2025.}.

In altre parole, tutta la complessità di COMPAS ,i suoi 137 fattori, i suoi algoritmi proprietari, la sua aura di scientificità ,si riduceva essenzialmente a dire: "I giovani con molti precedenti penali tendono a commettere altri reati". Una conclusione che qualsiasi persona di buon senso avrebbe potuto raggiungere senza bisogno di un algoritmo da milioni di dollari.

\section{Quando la tecnologia incontra i bias: una prospettiva europea}

L'Europa ha adottato un approccio più cauto all'uso dell'intelligenza artificiale nel sistema giudiziario. Nel 2018, la Commissione Europea per l'Efficienza della Giustizia (CEPEJ) ha adottato la "Carta etica sull'uso dell'intelligenza artificiale nei sistemi giudiziari", stabilendo principi fondamentali come la trasparenza, l'equità e la prevenzione delle discriminazioni\footnote{Agenda Digitale, "AI e giusto processo, facciamo il punto: le norme, le applicazioni, le sentenze", 2024.}.

La Carta distingue chiaramente tra ambito civile, commerciale e amministrativo ,dove l'IA può contribuire a migliorare prevedibilità e coerenza delle decisioni ,e ambito penale, dove "il suo utilizzo deve essere esaminato con le massime riserve, al fine di prevenire discriminazioni basate su dati sensibili".

È un approccio che riconosce quello che abbiamo imparato nel corso di questo libro: la tecnologia non è neutrale, e quando amplifica i nostri bias cognitivi in contesti ad alto rischio come la giustizia penale, le conseguenze possono essere catastrofiche.

In Italia, il Consiglio di Stato ha iniziato a delineare i principi per l'uso di algoritmi nella pubblica amministrazione. In una sentenza del 2019, i giudici amministrativi hanno stabilito che l'utilizzo di algoritmi deve garantire trasparenza, verificabilità e rispetto dei principi di ragionevolezza e proporzionalità\footnote{Consiglio di Stato, Sezione VI, Sentenza n. 2270 dell'8 aprile 2019.}.

\subsection{La ricerca di soluzioni: tra consapevolezza e design istituzionale}

Come abbiamo visto nei capitoli precedenti, la consapevolezza dei bias non li elimina magicamente. Anche Kahneman, dopo decenni di studio sui bias cognitivi, continua a caderci. Ma la consapevolezza è il primo passo per sviluppare strategie di mitigazione.

Nel campo giudiziario, queste strategie operano su più livelli. Alcuni sistemi giudiziari stanno sperimentando con procedure "cieche" simili a quelle utilizzate nella ricerca medica: analisi forensi condotte senza sapere quale campione appartiene al sospettato, testimonianze raccolte senza informazioni sul background dell'imputato, decisioni prese in base a trascrizioni anonimizzate.

Altri paesi stanno investendo nella formazione specifica sui bias cognitivi per magistrati e operatori della giustizia. Uno studio della Washington State University del 2023 ha dimostrato che programmi di formazione ben strutturati ,che includevano simulazioni realistiche e analisi dei propri comportamenti ,potevano ridurre significativamente le discriminazioni nelle decisioni delle forze dell'ordine\footnote{Washington State University, "Test of police implicit bias training shows modest improvement", 2023.}.

Ma forse l'approccio più promettente è quello del "design istituzionale": strutturare deliberatamente i processi decisionali per ridurre l'impatto dei bias. Ad esempio, richiedere che certe decisioni siano prese da gruppi diversificati piuttosto che da singoli individui, introdurre pause obbligatorie prima di decisioni importanti, standardizzare i formati di presentazione delle prove per ridurre effetti di framing.

\subsection{La lezione per tutti noi}

La storia dei bias cognitivi nel sistema giudiziario ci offre una lezione che va ben oltre le aule di tribunale. È la dimostrazione più chiara di come i meccanismi che abbiamo ereditato dal Pleistocene ,il pensiero veloce del Sistema 1, la ricerca di conferme, l'appartenenza al gruppo, l'ancoraggio a prime impressioni ,continuino a operare anche nei contesti più formali e regolamentati della nostra società.

Se persino giudici addestrati, vincolati da procedure rigorose e consapevoli dell'importanza delle loro decisioni cadono sistematicamente negli stessi tranelli cognitivi che affliggono tutti noi, questo ci dice qualcosa di fondamentale sulla natura umana: i bias non sono bug da correggere, ma feature evolutive così profondamente radicate che eliminarle significherebbe smettere di essere umani.

Ma questo non significa arrendersi. Come abbiamo visto, la combinazione di consapevolezza scientifica, design istituzionale intelligente e tecnologie ben progettate può aiutarci a costruire sistemi più equi. Non elimineremo mai completamente i bias ,né sarebbe desiderabile farlo ,ma possiamo imparare a riconoscerli, a mitigarne gli effetti più dannosi, e a costruire processi che siano più resistenti alle nostre limitazioni cognitive.

In fondo, è questa la sfida centrale del nostro tempo: come essere cavernicoli con lo smartphone senza che lo smartphone amplifichi i nostri istinti più tribali, e senza che i nostri istinti più tribali sabotino le potenzialità dello smartphone. Nel sistema giudiziario, come negli investimenti, nella medicina e in tutti gli altri ambiti che abbiamo esplorato, la risposta non sta nell'eliminare la componente umana, ma nel progettare sistemi che portino fuori il meglio della nostra umanità, contenendo il peggio.

Perché in fondo, anche se il nostro cervello del Pleistocene spesso ci gioca brutti scherzi, è lo stesso cervello che ci ha permesso di sviluppare la scienza per studiarci, la tecnologia per potenziarci, e la saggezza per riconoscere i nostri limiti. E questa, forse, è la cosa più umana di tutte.
